\documentclass[12pt,a4paper]{article}
\usepackage{amsmath,amssymb,amsthm,hyperref,geometry}
\usepackage{enumitem,booktabs,array}
\usepackage[T1]{fontenc}
\usepackage[utf8]{inputenc}
\usepackage[ngerman]{babel}
\geometry{margin=2.5cm}

% -------------------------------------------------------
% Theorem-Umgebungen
% -------------------------------------------------------
\newtheorem{theorem}{Satz}[section]
\newtheorem{lemma}[theorem]{Lemma}
\newtheorem{corollary}[theorem]{Korollar}
\newtheorem{proposition}[theorem]{Proposition}
\newtheorem{conjecture}[theorem]{Vermutung}
\newtheorem{definition}[theorem]{Definition}

\theoremstyle{remark}
\newtheorem{remark}[theorem]{Bemerkung}
\newtheorem{example}[theorem]{Beispiel}

% -------------------------------------------------------
% Metadaten
% -------------------------------------------------------
\title{\textbf{Navier--Stokes-Gleichungen: Existenz und Regularität}\\[0.5em]
\large Schwache Lösungen, Regularität, Turbulenz und das Millennium-Problem}
\author{Michael Fuhrmann}
\date{März 2026}

\hypersetup{
  pdftitle={Navier-Stokes-Existenz und Regularität},
  pdfauthor={Michael Fuhrmann},
  colorlinks=true,
  linkcolor=blue,
  citecolor=blue
}

\begin{document}
\maketitle
\begin{abstract}
Die Navier--Stokes-Gleichungen beschreiben die Bewegung zäher inkompressibler Flüssigkeiten
und bilden die Grundlage nahezu aller Modellierungen in der Strömungslehre, von der
Aeronautik bis zur Wettervorhersage. Das Millennium-Preisproblem fragt, ob für beliebige
glatte, schnell abfallende Anfangsdaten $u_0 \in C^\infty(\mathbb{R}^3)$ immer glatte
(unendlich oft differenzierbare) globale Lösungen existieren, oder ob Singularitäten
(„Blow-up") in endlicher Zeit entstehen können.

Dieses Paper stellt die Gleichungen, die benötigten Funktionenräume (Sobolev, Leray) für
schwache Lösungen, Lerays fundamentale Ergebnisse von 1934, Ladyzhenskayas vollständige
2D-Theorie, den Partialregularitätssatz von Caffarelli--Kohn--Nirenberg, Kolmogorovs
Turbulenztheorie von 1941 sowie Taos gemittelte Navier--Stokes-Blow-up-Konstruktion von
2016 vor. Der aktuelle Stand im Jahr 2026 wird zusammengefasst.
\end{abstract}

\tableofcontents
\newpage

% ================================================================
\section{Die Navier--Stokes-Gleichungen}
% ================================================================

\subsection{Herleitung und physikalische Bedeutung}

Die \emph{Navier--Stokes-Gleichungen} für eine inkompressible zähe Newtonsche Flüssigkeit
mit Geschwindigkeitsfeld $u\colon \mathbb{R}^3 \times [0,\infty) \to \mathbb{R}^3$ und
Druck $p\colon \mathbb{R}^3 \times [0,\infty) \to \mathbb{R}$ lauten:
\begin{align}
  \partial_t u + (u \cdot \nabla)u &= -\nabla p + \nu \Delta u + f, \label{eq:ns_impuls}\\
  \nabla \cdot u &= 0, \label{eq:ns_inkompressibel}
\end{align}
mit Anfangsbedingung $u(x, 0) = u_0(x)$.

Dabei bedeuten:
\begin{itemize}[nosep]
  \item $(u \cdot \nabla)u$: der \emph{Advektionsterm} (nichtlinear),
  \item $-\nabla p$: der \emph{Druckgradient},
  \item $\nu > 0$: die \emph{kinematische Viskosität},
  \item $\nu \Delta u$: der \emph{Diffusionsterm} (Viskositätsspannung),
  \item $f$: eine äußere Volumenkraft (z.\,B.\ Schwerkraft).
\end{itemize}

Gleichung \eqref{eq:ns_inkompressibel} drückt die Inkompressibilitätsbedingung
$\operatorname{div}\, u = 0$ aus, die besagt, dass die Flüssigkeit volumenerhaltend ist.

\subsection{Die Millennium-Problemstellung}

Das Clay Mathematics Institute~\cite{Clay2000} verlangt einen der folgenden Beweise:
\begin{enumerate}[nosep]
  \item Einen Beweis, dass für beliebige glatte, schnell abfallende Anfangsdaten
    $u_0 \in C^\infty(\mathbb{R}^3)$ mit $\nabla \cdot u_0 = 0$ immer eine glatte globale
    Lösung $(u, p)$ mit $u \in C^\infty(\mathbb{R}^3 \times [0,\infty))$ existiert.
  \item Einen Beweis eines \emph{Blow-ups}: Anfangsdaten, für die keine glatte globale
    Lösung existiert.
\end{enumerate}

Das Problem ist in $\mathbb{R}^3$ (und auf dem Torus $\mathbb{T}^3$) offen.

% ================================================================
\section{Funktionenräume und schwache Formulierung}
% ================================================================

\subsection{Sobolev-Räume}

\begin{definition}[Sobolev-Raum $H^k$]
  Für $k \geq 0$ besteht der Sobolev-Raum $H^k(\mathbb{R}^n) = W^{k,2}(\mathbb{R}^n)$
  aus Funktionen $f \in L^2(\mathbb{R}^n)$ mit schwachen Ableitungen bis zur Ordnung $k$
  in $L^2$:
  \[
    H^k(\mathbb{R}^n) \;=\; \bigl\{ f \in L^2(\mathbb{R}^n) :
    \|\hat{f}\|_{H^k}^2 := \int_{\mathbb{R}^n} (1+|\xi|^2)^k |\hat{f}(\xi)|^2\,d\xi < \infty \bigr\}.
  \]
\end{definition}

Die Energie der Navier--Stokes-Lösung wird durch die $L^2$-Norm von $u$ gemessen:
\[
  E(t) \;=\; \frac{1}{2}\int_{\mathbb{R}^3} |u(x,t)|^2\,dx.
\]
Formal ergibt Multiplikation von \eqref{eq:ns_impuls} mit $u$ und Integration die
\emph{Energieungleichung}:
\begin{equation}\label{eq:energie}
  E(t) + \nu \int_0^t \int_{\mathbb{R}^3} |\nabla u|^2\,dx\,ds \;\leq\; E(0),
\end{equation}
für $f \equiv 0$. Dies zeigt, dass Energie nicht zunimmt (Viskosität dissipiert Energie).

\subsection{Lerays schwache Formulierung (1934)}

\begin{definition}[Schwache (Leray-) Lösung~\cite{Leray1934}]
  Eine \emph{schwache Lösung} auf $[0,T]$ ist ein divergenzfreies Vektorfeld
  $u \in L^\infty(0,T; L^2) \cap L^2(0,T; H^1)$, das folgende Bedingungen erfüllt:
  \begin{enumerate}[nosep]
    \item Für alle divergenzfreien Testfunktionen $\phi \in C_c^\infty(\mathbb{R}^3 \times (0,T))$:
    \[
      \int_0^T\!\int_{\mathbb{R}^3} \!\bigl(-u \cdot \partial_t \phi
      - (u\otimes u):\nabla\phi + \nu \nabla u : \nabla\phi\bigr)\,dx\,dt = 0.
    \]
    \item Die Energieungleichung \eqref{eq:energie} gilt.
    \item $u(t) \to u_0$ in $L^2$ für $t \to 0^+$.
  \end{enumerate}
\end{definition}

\begin{theorem}[Leray 1934~\cite{Leray1934}]
  Für jedes $u_0 \in L^2(\mathbb{R}^3)$ mit $\nabla \cdot u_0 = 0$ und $f \in L^2$
  existiert mindestens eine globale schwache (Leray-) Lösung.
\end{theorem}

Lerays Beweis verwendet eine \emph{Regularisierung} (Glättung der Nichtlinearität durch
einen Parameter $\varepsilon$) und anschließend den Grenzübergang $\varepsilon \to 0$
unter Verwendung von Kompaktheit (Aubin--Lions-Lemma). Die entstehende Lösung wird
\emph{Leray--Hopf-schwache Lösung} nach Leray~\cite{Leray1934} und Hopf~\cite{Hopf1951}
genannt.

% ================================================================
\section{Der 2D-Fall: Vollständige Theorie}
% ================================================================

\begin{theorem}[Ladyzhenskaya 1969~\cite{Lady1969}]
  In zwei Raumdimensionen ($u\colon \mathbb{R}^2 \times [0,\infty) \to \mathbb{R}^2$)
  gilt für jedes $u_0 \in H^1(\mathbb{R}^2)$ mit $\nabla \cdot u_0 = 0$:
  \begin{enumerate}[nosep]
    \item Es existiert eine eindeutige globale glatte Lösung.
    \item Die Lösung hängt stetig von den Anfangsdaten ab.
    \item Die Strömung ist global wohlgestellt in $H^k$ für alle $k \geq 1$.
  \end{enumerate}
\end{theorem}

Das zentrale Hilfsmittel in 2D ist die \emph{Ladyzhenskaya-Ungleichung}:
\[
  \|u\|_{L^4}^2 \;\leq\; C\|u\|_{L^2}\|\nabla u\|_{L^2},
\]
die zusammen mit der 2D-Wirbelgleichung $\partial_t \omega + (u\cdot\nabla)\omega = \nu\Delta\omega$
(wobei der Wirbelstreck-Term $\omega\cdot\nabla u$ in 2D verschwindet) eine $L^\infty$-Schranke
für die Wirbelstärke liefert und einen Blow-up verhindert.

In 3D lautet die entsprechende Wirbelgleichung:
$\partial_t \omega + (u\cdot\nabla)\omega = (\omega\cdot\nabla)u + \nu\Delta\omega$,
wobei der \emph{Wirbelstreck-Term} $(\omega\cdot\nabla)u$ die Wirbelstärke verstärken und
potenziell zu einem Blow-up führen kann.

% ================================================================
\section{Das 3D-Regularitätsproblem}
% ================================================================

\subsection{Lokale Existenz glatter Lösungen}

\begin{theorem}[Kato--Fujita 1962~\cite{KatoFujita1962}]
  Für $u_0 \in H^s(\mathbb{R}^3)$ mit $s \geq 1/2$ und $\nabla\cdot u_0=0$ existieren
  $T^* > 0$ und eine eindeutige glatte Lösung $u \in C([0,T^*); H^s)$.
\end{theorem}

Die Frage ist, ob die maximale Existenzzeit $T^*$ immer $+\infty$ ist (globale Existenz)
oder endlich sein kann (Blow-up).

\subsection{Blow-up-Kriterien}

\begin{theorem}[Beale--Kato--Majda 1984~\cite{BKM1984}]
  Falls $T^* < \infty$ (Blow-up in endlicher Zeit), muss notwendigerweise
  \[
    \int_0^{T^*} \|\omega(\cdot,t)\|_{L^\infty}\,dt = +\infty
  \]
  gelten, wobei $\omega = \nabla \times u$ die Wirbelstärke ist. Äquivalent: Die
  Wirbelstärke muss für einen Blow-up unbeschränkt werden.
\end{theorem}

Dieses Kriterium lokalisiert das Regularitätsproblem auf das Verhalten der Wirbelstärke.

\subsection{Serrins Regularitätskriterium}

\begin{theorem}[Serrin 1962~\cite{Serrin1962}; Prodi 1959]
  Sei $u$ eine Leray--Hopf-schwache Lösung auf $(0,T)$. Falls zusätzlich
  $u \in L^q(0,T; L^r(\mathbb{R}^3))$ für ein Paar $(q,r)$ mit
  \[
    \frac{2}{q} + \frac{3}{r} = 1, \quad r > 3,
  \]
  gilt, so ist $u$ die eindeutige glatte Lösung auf $(0,T)$.
\end{theorem}

Die Serrin-Klassen-Schranken sind skaleninvariant (die Navier--Stokes-Gleichungen sind
invariant unter $(u,p,x,t) \mapsto (\lambda u, \lambda^2 p, \lambda^{-1}x, \lambda^{-2}t)$),
und der kritische Fall $r=3$, $q=\infty$ wurde durch Escauriaza--Seregin--Šverák
(2003)~\cite{ESS2003} bewiesen.

% ================================================================
\section{Caffarelli--Kohn--Nirenberg-Partialregularität}
% ================================================================

\begin{theorem}[Caffarelli--Kohn--Nirenberg 1982~\cite{CKN1982}]
  Sei $(u, p)$ eine \emph{geeignete schwache Lösung} auf $\mathbb{R}^3 \times (0,T)$
  (eine Leray--Hopf-Lösung, die eine lokale Energieungleichung erfüllt). Dann genügt die
  singuläre Menge $\mathcal{S} = \{(x,t) : u \text{ ist lokal unbeschränkt bei } (x,t)\}$
  der Abschätzung:
  \[
    \mathcal{H}^1(\mathcal{S}) = 0,
  \]
  wobei $\mathcal{H}^1$ das eindimensionale parabolische Hausdorff-Maß ist.
  Insbesondere hat die singuläre Menge \emph{parabolische Hausdorff-Dimension} $\leq 1$.
\end{theorem}

Dies ist das stärkste unbedingte Regularitätsergebnis in 3D. Insbesondere:
\begin{itemize}[nosep]
  \item Die singuläre Menge kann keine Kurve, Fläche oder offene Menge enthalten.
  \item Zu jedem festen Zeitpunkt $t$ kann es höchstens endlich viele singuläre Punkte geben.
  \item Raumzeit-Singularitäten (falls vorhanden) sind sehr dünn gesät.
\end{itemize}

\begin{remark}
  Der CKN-Satz schließt nicht aus, dass $\mathcal{S}$ nicht-leer ist; er schränkt lediglich
  ihre Größe ein. Die Konstruktion einer geeigneten schwachen Lösung mit nicht-leerer
  singulärer Menge, oder der Beweis $\mathcal{S} = \emptyset$ für alle glatten Daten,
  bleibt offen.
\end{remark}

% ================================================================
\section{Turbulenz: Kolmogorovs Theorie von 1941}
% ================================================================

\subsection{Die Energiekaskade}

Kolmogorov~\cite{Kolmogorov1941} (1941) schlug eine statistische Theorie voll entwickelter
Turbulenz auf der Grundlage zweier Hypothesen vor:
\begin{enumerate}[nosep]
  \item \textbf{Lokale Isotropie:} Bei hohen Reynolds-Zahlen sind kleinskalige turbulente
    Bewegungen statistisch isotrop (unabhängig von der Großskalen-Geometrie).
  \item \textbf{Universelles Gleichgewicht:} Im \emph{Inertialbereich} (Skalen zwischen
    der Integralskala $L$ und der viskosen Dissipationsskala $\eta$) hängt die Statistik
    nur von der mittleren Energiedissipationsrate $\varepsilon$ und der Viskosität $\nu$ ab.
\end{enumerate}

\subsection{Das Kolmogorov-$-5/3$-Spektrum}

Aus dimensionsanalytischen Gründen allein (mit $[\varepsilon] = \text{m}^2\text{s}^{-3}$
und $[k] = \text{m}^{-1}$) muss das Energiespektrum $E(k)$ (Energie pro Wellenzahleinheit)
der Beziehung genügen:
\begin{equation}\label{eq:k41}
  E(k) \;=\; C_K \varepsilon^{2/3} k^{-5/3}, \quad \eta^{-1} \ll k \ll L^{-1},
\end{equation}
wobei $C_K \approx 1{,}5$ die \emph{Kolmogorov-Konstante} ist.

Die Kolmogorov-Dissipationsskala ist $\eta = (\nu^3/\varepsilon)^{1/4}$ und die Reynolds-Zahl
$\mathrm{Re} = L^{4/3}\varepsilon^{1/3}/\nu$. Das $k^{-5/3}$-Gesetz \eqref{eq:k41}
ist experimentell über viele Größenordnungen bestätigt.

\subsection{K41 und Navier--Stokes}

Eine rigorose Herleitung von K41 aus den Navier--Stokes-Gleichungen ist bislang nicht
gelungen. Constantin und Doering~\cite{ConstantinDoering1994} haben obere Schranken für
die Energiedissipation bewiesen, die mit K41 verträglich sind. Die Verbindung zwischen der
K41-Phänomenologie und dem mathematischen Regularitätsproblem (Blow-up) ist ein aktives
Forschungsgebiet.

% ================================================================
\section{Taos gemittelte Navier--Stokes-Blow-up-Konstruktion (2016)}
% ================================================================

\subsection{Die gemittelten Gleichungen}

Im Jahr 2016 konstruierte Terence Tao~\cite{Tao2016} einen Blow-up für ein modifiziertes
System, die sogenannten \emph{gemittelten Navier--Stokes-Gleichungen}:
\[
  \partial_t u = \nu\Delta u - B(u,u), \quad \nabla\cdot u = 0,
\]
wobei $B(u,v) = \frac{1}{2}[\nabla\times(u\times v) + P\nabla(u\cdot v)]$ eine
symmetrisierte Bilinearform ist, die mit der Navier--Stokes-Nichtlinearität verwandt und
über SO(3)-Rotationen gemittelt wurde:
\[
  B_{\mathrm{avg}}(u,v) \;=\; \mathbb{E}_{R \in SO(3)}\bigl[R^{-1}B(Ru, Rv)\bigr].
\]

\subsection{Taos Satz}

\begin{theorem}[Tao 2016~\cite{Tao2016}]
  Es gibt ein gemitteltes Navier--Stokes-System mit einer glatten, divergenzfreien
  Anfangsbedingung $u_0 \in C^\infty(\mathbb{R}^3)$, sodass die eindeutige glatte Lösung
  $u(t)$ zu einem endlichen Zeitpunkt $T^* < \infty$ aufbläst:
  \[
    \limsup_{t \to T^*} \|u(t)\|_{H^1(\mathbb{R}^3)} = +\infty.
  \]
\end{theorem}

\begin{remark}
  Das gemittelte System erfüllt die \emph{Energiegleichung} (nicht nur die Ungleichung),
  ist divergenzfrei und hat dieselbe Skaleninvarianz wie Navier--Stokes. Es unterscheidet
  sich von den wahren Navier--Stokes-Gleichungen darin, dass die Nichtlinearität gemittelt
  ist.
\end{remark}

\subsection{Implikationen von Taos Ergebnis}

Taos Paper~\cite{Tao2016} erfüllt mehrere Zwecke:
\begin{enumerate}[nosep]
  \item Es zeigt, dass ein Blow-up durch die Struktur der Gleichungen allein
    (Energieungleichung, Divergenzfreiheit, Skalierung) \emph{nicht ausgeschlossen} wird.
  \item Der Blow-up-Mechanismus (eine „Energie-Kaskaden-Maschine") liefert ein
    konzeptionelles Modell dafür, wie Singularitäten im wahren Navier--Stokes-System
    entstehen könnten.
  \item Er identifiziert Hindernisse: Jeder Beweis globaler Regularität muss eine
    Eigenschaft von Navier--Stokes nutzen, die das gemittelte System nicht besitzt.
\end{enumerate}

% ================================================================
\section{Neuere Entwicklungen (2020--2026)}
% ================================================================

\begin{itemize}[nosep]
  \item \textbf{Nicht-Eindeutigkeit schwacher Lösungen (2019--2022):} Buckmaster und
    Vicol~\cite{BuckmasterVicol2019} bewiesen mittels konvexer Integration, dass
    nicht-eindeutige schwache Lösungen mit $H^{1/3-\varepsilon}$-Regularität existieren.
    Dies legt nahe, dass Leray-Lösungen unterhalb einer kritischen Regularitätsschwelle
    nicht eindeutig sind.
  \item \textbf{Onsager-Vermutung (bewiesen):} De Lellis und Székelyhidi (2019)~\cite{DeLellis2019}
    bewiesen die Onsager-Vermutung: Es existieren Lösungen der \emph{Euler-Gleichungen}
    ($\nu=0$) mit $C^{1/3-\varepsilon}$-Regularität, die Energie dissipieren, während alle
    Lösungen mit $C^{1/3+\varepsilon}$-Regularität Energie erhalten.
  \item \textbf{Stochastische Navier--Stokes:} Das Hinzufügen von Rauschen kann
    Eindeutigkeit wiederherstellen (Da Prato et al.); Eindeutigkeit für stochastische
    Navier--Stokes in 3D in bestimmten Klassen wurde von Hofmanova et al.\ (2022)
    nachgewiesen.
  \item \textbf{Maschinelles Lernen und Turbulenz (2023--2026):} Neuronale Operatoren
    (Fourier-Neural-Operators von Li et al.) haben bemerkenswerte Genauigkeit bei
    Turbulenzsimulationen erzielt, sind aber numerisch-statistische Ansätze ohne Bezug
    zum mathematischen Existenzproblem.
  \item \textbf{Millennium-Problem (2026):} Kein Beweis globaler Regularität oder eines
    Blow-ups für die wahren 3D-Navier--Stokes-Gleichungen wurde gefunden. Das Problem
    bleibt vollständig offen.
\end{itemize}

% ================================================================
\section{Zusammenfassung der Ergebnisse}
% ================================================================

\begin{table}[h]
\centering
\begin{tabular}{lll}
\toprule
Dimension & Existenz & Eindeutigkeit/Regularität \\
\midrule
2D & Global, glatt (Leray 1934) & Eindeutig, glatt (Ladyzhenskaya 1969) \\
3D & Global schwach (Leray 1934) & Offen \\
3D (kurze Zeit) & Lokal glatt (Kato--Fujita 1962) & Eindeutig (kurze Zeit) \\
3D (sing.\ Menge) & Hausdorff-Dim.\ $\leq 1$ (CKN 1982) & --- \\
Gemittelt 3D & Blow-up (Tao 2016) & --- \\
\bottomrule
\end{tabular}
\caption{Überblick über die Navier--Stokes-Ergebnisse}
\end{table}

% ================================================================
\section{Schlussfolgerung}
% ================================================================

Das Navier--Stokes-Problem ist ein Paradebeispiel für die tiefe Kluft zwischen
physikalischer Intuition und mathematischer Strenge. Flüssigkeiten verhalten sich in jedem
physikalischen Experiment und jeder numerischen Simulation glatt, doch dieses Verhalten
aus den Gleichungen heraus zu beweisen, hat sich fast ein Jahrhundert lang allen
Bemühungen widersetzt.

Die Teilergebnisse — Lerays schwache Lösungen, Ladyzhenskayas vollständige 2D-Theorie,
Caffarelli--Kohn--Nirenbergs Singularitätsmengen-Schranke und Taos gemittelter Blow-up —
umreißen zusammen die Grenzen des aktuellen Wissens. Die Lösung des Millennium-Problems
wird wahrscheinlich echte neue Mathematik erfordern: neue Funktionenräume, neue
Regularitätskriterien oder eine unerwartete Blow-up-Konstruktion.

\begin{thebibliography}{99}

\bibitem{Clay2000}
C.~L.~Fefferman,
\textit{Existence and Smoothness of the Navier--Stokes Equation},
Beschreibung des Clay Mathematics Institute Millennium-Problems, 2000,
\url{https://www.claymath.org/millennium-problems/navier-stokes-equation}.

\bibitem{Leray1934}
J.~Leray,
\textit{Sur le mouvement d'un liquide visqueux emplissant l'espace},
Acta Math.\ \textbf{63} (1934), 193--248.

\bibitem{Hopf1951}
E.~Hopf,
\textit{Über die Anfangswertaufgabe für die hydrodynamischen Grundgleichungen},
Math.\ Nachr.\ \textbf{4} (1951), 213--231.

\bibitem{Lady1969}
O.~A.~Ladyzhenskaya,
\textit{The Mathematical Theory of Viscous Incompressible Flow},
Gordon and Breach, New York, 1969.

\bibitem{KatoFujita1962}
T.~Kato und H.~Fujita,
\textit{On the non-stationary Navier--Stokes system},
Rend.\ Sem.\ Mat.\ Univ.\ Padova \textbf{32} (1962), 243--260.

\bibitem{BKM1984}
J.~T.~Beale, T.~Kato und A.~Majda,
\textit{Remarks on the breakdown of smooth solutions for the 3-D Euler equations},
Comm.\ Math.\ Phys.\ \textbf{94} (1984), 61--66.

\bibitem{Serrin1962}
J.~Serrin,
\textit{On the interior regularity of weak solutions of the Navier--Stokes equations},
Arch.\ Rational Mech.\ Anal.\ \textbf{9} (1962), 187--195.

\bibitem{ESS2003}
L.~Escauriaza, G.~Seregin und V.~Šverák,
\textit{$L_{3,\infty}$-solutions of the Navier--Stokes equations and backward uniqueness},
Russian Math.\ Surveys \textbf{58} (2003), 211--250.

\bibitem{CKN1982}
L.~Caffarelli, R.~Kohn und L.~Nirenberg,
\textit{Partial regularity of suitable weak solutions of the Navier--Stokes equations},
Comm.\ Pure Appl.\ Math.\ \textbf{35} (1982), 771--831.

\bibitem{Kolmogorov1941}
A.~N.~Kolmogorov,
\textit{The local structure of turbulence in incompressible viscous fluid for very large Reynolds numbers},
Dokl.\ Akad.\ Nauk SSSR \textbf{30} (1941), 301--305.

\bibitem{ConstantinDoering1994}
P.~Constantin und C.~R.~Doering,
\textit{Variational bounds on energy dissipation in incompressible flows: II. Channel flow},
Phys.\ Rev.\ E \textbf{51} (1995), 3192--3198.

\bibitem{Tao2016}
T.~Tao,
\textit{Finite time blowup for an averaged three-dimensional Navier--Stokes equation},
J.\ Amer.\ Math.\ Soc.\ \textbf{29} (2016), 601--674.

\bibitem{BuckmasterVicol2019}
T.~Buckmaster und V.~Vicol,
\textit{Nonuniqueness of weak solutions to the Navier--Stokes equation},
Ann.\ of Math.\ \textbf{189} (2019), 101--144.

\bibitem{DeLellis2019}
C.~De Lellis und L.~Székelyhidi,
\textit{On turbulence and geometry: from Nash to Onsager},
Notices Amer.\ Math.\ Soc.\ \textbf{66} (2019), 677--685.

\bibitem{ConstantinFoias1988}
P.~Constantin und C.~Foias,
\textit{Navier--Stokes Equations},
University of Chicago Press, 1988.

\bibitem{Temam1984}
R.~Temam,
\textit{Navier--Stokes Equations: Theory and Numerical Analysis},
AMS Chelsea Publishing, 1984.

\end{thebibliography}

\end{document}
